%Esa en la hipótesis de partida a partir de la cual compuse mi propia relectura de Ulises, en \emph{Abandonando Itaca}\footnote{https://eolasediciones.es/libro/abandonando-itaca-leaving-ithaca/}. Parte de mis razones para hacerlo tienen que ver con mi propia \emph{Odisea}, ese viaje por la vida en el que todos andamos buscando Ítaca, olvidando a menudo los versos de Kavafis\footnote{En traducción de J.M. García, ed. Hiperión}:
%
%Si vas a emprender el viaje hacia Ítaca,
%pide que tu camino sea largo,
%rico en experiencias, en conocimiento.
%
%La primera vez que escuché estos versos, acababa de cumplir los quince años. Estaba pasando el verano en la Costa Brava, en Blanes, con la familia de mi amigo Quelo\footnote{Necesitaría un poemario completo ---en realidad lo tengo--- para hablar de Miguel Asensi, mi amigohermano, mi hermanoamigo. La canción de Serrat, <<decir amigo>>, se aproxima bastante a lo que éramos. Murió hace cinco años, de un cáncer mal diagnosticado, en la época de la pandemia. Acababa de cumplir los sesenta.} y los dos trabajábamos para la Siseta, que tenía un lucrativo negocio, vendiéndoles fotos de recuerdo a los turistas. Hay un relato corto, <<El otro Miguel Strogoff>>, recién publicado\footnote{https://www.todoliteratura.es/noticia/63055/criticas/el-dilema-turing-y-otros-cuentos-de-juan-jose-gomez-cadenas-ilustraciones-de-miguel-almagro-juan-manuel-fontenls-juan-gomez-macias-rafael-leonardo-setien-francis-marchena-luis-ruiz-padron.html}, que captura las experiencias de aquel verano. En aquellos meses de julio y agosto, pasaron dos cosas importantes en, mi entonces, corta vida. Sobreviví, como el correo del zar al hierro candente ---del desengaño--- y Lluis Llach me descubrió los versos de Kavafis.
%
%Quan surts per fer el viatge cap a Itaca,
%has de pregar que el camí sigui llarg,
%ple d'aventures, ple de coneixences.
%
%No sabía yo entonces que los versos reproducían fielmente al griego, pero la música de aquel álbum nos hacía creer que las golondrinas atestadas de turistas en las que nos pasábamos el día ---vendiendo fotos y souvenirs--- eran las cóncavas naves de los aqueos.
%
%Muchas veces he pensado que mi nave zarpó de Ítaca ese verano, dejando atrás mi niñez, llevándome por océanos que ya llevo, desde entonces, medio siglo transitando. En algún momento de ese viaje, hace ya más de treinta y cinco años, arribé a la isla de la hechicera de las largas trenzas. Y ya no pude, o no quise, volver a Ítaca.
%
%Tanto el Ulises del texto clásico como el protagonista de Nolan llegan a